
% --------------------------------------------------------------------------

\section{Plan de trabajo}
\label{sc:PT}
La planificación esta desglosada en 2 paquetes de trabajo (WP). El primero (WP$_{1}$) corresponde a la implementación de la técnica híbrida en las metaheurísticas seleccionadas con pruebas preliminares de su funcionamiento. Luego, el segundo paquete (WP$_{2}$) tratará la fase de experimentación junto al análisis de resultados, con la posterior interpretación de estos y conclusión. El trabajo de paquetes (WP) están separados en actividades (A$_{x.y}$) para presentar el plan de trabajo mejor (ver la Tabla~\ref{tbl:wp}).



\begin{comment}

\begin{table}[!ht]
\begin{center}
\begin{scriptsize}
\caption{Gráfico Gantt - Semestre 2, 2024}\label{tbl:wp}
\vspace{5pt}
\begin{tabular}{|p{.42cm}|p{.15cm}|p{.15cm}|p{.15cm}|p{.15cm}|p{.15cm}|p{.15cm}|p{.15cm}|p{.15cm}|p{.15cm}|p{.15cm}|p{.15cm}|p{.15cm}|p{.15cm}|p{.15cm}|p{.15cm}|p{.15cm}|p{.15cm}|p{.15cm}|p{.15cm}|p{.15cm}|}
\hline
\rowcolor{gray!20}
\multicolumn{1}{|c|}{}
& \multicolumn{4}{c|}{\textbf{AGOSTO}}
& \multicolumn{4}{c|}{\textbf{SEPTIEMBRE}}
& \multicolumn{4}{c|}{\textbf{OCTUBRE}}
& \multicolumn{4}{c|}{\textbf{NOVIEMBRE}}
& \multicolumn{4}{c|}{\textbf{DICIEMBRE}} \\ \cline{2-21}
\rowcolor{gray!20}
\multicolumn{1}{|c|}{\multirow{-2}{*}{\textbf{Act.}}} &
\multicolumn{1}{c|}{\textbf{1}} &
\multicolumn{1}{c|}{\textbf{2}} &
\multicolumn{1}{c|}{\textbf{3}} &
\multicolumn{1}{c|}{\textbf{4}} &
\multicolumn{1}{c|}{\textbf{1}} &
\multicolumn{1}{c|}{\textbf{2}} &
\multicolumn{1}{c|}{\textbf{3}} &
\multicolumn{1}{c|}{\textbf{4}} &
\multicolumn{1}{c|}{\textbf{1}} &
\multicolumn{1}{c|}{\textbf{2}} &
\multicolumn{1}{c|}{\textbf{3}} &
\multicolumn{1}{c|}{\textbf{4}} &
\multicolumn{1}{c|}{\textbf{1}} &
\multicolumn{1}{c|}{\textbf{2}} &
\multicolumn{1}{c|}{\textbf{3}} &
\multicolumn{1}{c|}{\textbf{4}} &
\multicolumn{1}{c|}{\textbf{1}} &
\multicolumn{1}{c|}{\textbf{2}} &
\multicolumn{1}{c|}{\textbf{3}} &
\multicolumn{1}{c|}{\textbf{4}} \\ \hline
\multicolumn{1}{c}{\textbf{WP}$_{1}$} \\ \hline
A$_{1.1}$ &X&X&X&X& & & & & & & & & & & & & & & & \\ \hline
A$_{1.2}$ & & &X&X&X&X&X& & & & & & & & & & & & & \\ \hline
A$_{1.3}$ & & &X&X& & & & & & & & & & & & & & & & \\ \hline
A$_{1.4}$ & & & & & & & &X&X&X&X&X&X& & & & & & & \\ \hline
A$_{1.5}$ & & & & & & & & & & & &X&X&X& & & & & & \\ \hline
\multicolumn{1}{c}{\textbf{WP}$_{2}$} \\ \hline
A$_{2.1}$ & & & & & & & & & & & & & &X&X& & & & & \\ \hline
A$_{2.2}$ & & & & & & & & & & & & & & &X&X& & & & \\ \hline
A$_{2.3}$ & & & & & & & & & & & & & & & & &X&X& & \\ \hline
A$_{2.4}$ & & & & & & & & & & & & & & & & & &X&X&X \\ \hline
\end{tabular}
\end{scriptsize}
\end{center}
\end{table}

\end{comment}


\begin{table}[!ht]
\begin{center}
\begin{scriptsize}
\caption{Gráfico Gantt - Semestre 2, 2024}\label{tbl:wp}
\vspace{5pt}
\begin{tabular}{|p{.42cm}|p{.6cm}|p{.6cm}|p{.6cm}|p{.6cm}|p{.6cm}|p{.6cm}|p{.6cm}|p{.6cm}|p{.6cm}|p{.6cm}|p{.6cm}|p{.6cm}|p{.6cm}|p{.6cm}|p{.6cm}|p{.6cm}|p{.6cm}|p{.6cm}|p{.6cm}|p{.6cm}|}
\hline
\rowcolor{gray!20}
\multicolumn{1}{|c|}{}
& \multicolumn{4}{c|}{\textbf{AGOSTO}}
& \multicolumn{4}{c|}{\textbf{SEPTIEMBRE}}
& \multicolumn{4}{c|}{\textbf{OCTUBRE}}
& \multicolumn{4}{c|}{\textbf{NOVIEMBRE}}
& \multicolumn{4}{c|}{\textbf{DICIEMBRE}} \\ \cline{2-21}
\rowcolor{gray!20}
\multicolumn{1}{|c|}{\multirow{-2}{*}{\textbf{Act.}}} &
\multicolumn{1}{c|}{\textbf{1}} &
\multicolumn{1}{c|}{\textbf{2}} &
\multicolumn{1}{c|}{\textbf{3}} &
\multicolumn{1}{c|}{\textbf{4}} &
\multicolumn{1}{c|}{\textbf{1}} &
\multicolumn{1}{c|}{\textbf{2}} &
\multicolumn{1}{c|}{\textbf{3}} &
\multicolumn{1}{c|}{\textbf{4}} &
\multicolumn{1}{c|}{\textbf{1}} &
\multicolumn{1}{c|}{\textbf{2}} &
\multicolumn{1}{c|}{\textbf{3}} &
\multicolumn{1}{c|}{\textbf{4}} &
\multicolumn{1}{c|}{\textbf{1}} &
\multicolumn{1}{c|}{\textbf{2}} &
\multicolumn{1}{c|}{\textbf{3}} &
\multicolumn{1}{c|}{\textbf{4}} &
\multicolumn{1}{c|}{\textbf{1}} &
\multicolumn{1}{c|}{\textbf{2}} &
\multicolumn{1}{c|}{\textbf{3}} &
\multicolumn{1}{c|}{\textbf{4}} \\ \hline
\multicolumn{1}{c}{\textbf{WP}$_{1}$} \\ \hline
A$_{1.1}$ &X&X&X&X& & & & & & & & & & & & & & & & \\ \hline
A$_{1.2}$ & & &X&X&X&X&X& & & & & & & & & & & & & \\ \hline
A$_{1.3}$ & & &X&X& & & & & & & & & & & & & & & & \\ \hline
A$_{1.4}$ & & & & & & & &X&X&X&X&X&X& & & & & & & \\ \hline
A$_{1.5}$ & & & & & & & & & & & &X&X&X& & & & & & \\ \hline
\multicolumn{1}{c}{\textbf{WP}$_{2}$} \\ \hline
A$_{2.1}$ & & & & & & & & & & & & & &X&X& & & & & \\ \hline
A$_{2.2}$ & & & & & & & & & & & & & & &X&X& & & & \\ \hline
A$_{2.3}$ & & & & & & & & & & & & & & & & &X&X& & \\ \hline
A$_{2.4}$ & & & & & & & & & & & & & & & & & &X&X&  \\ \hline
\end{tabular}
\end{scriptsize}
\end{center}
\end{table}




\begin{itemize}

    \item \textbf{WP$_1$: Implementación de la técnica híbrida en las metaheurísticas para los problemas de optimización y ajuste de parámetros.}

    \begin{itemize}

        \item \textbf{A$_{1.1}$} Implementación de al menos 4 metaheurísticas para resolver funciones benchmark.

        \item \textbf{A$_{1.2}$} Estudiar, diseñar, y documentar en código los procedimientos de la propuesta para integrarlas a la metaheurística.

        \item \textbf{A$_{1.3}$} Ajustar los parámetros de las metaheurísticas a través de las recomendaciones hechas en los trabajos originales del algoritmo de optimización.

        \item \textbf{A$_{1.4}$} Implementar las metaheurísticas adaptadas al problema de optimización Profitable Tour Problem.

        \item \textbf{A$_{1.5}$} Diseñar e implementar métodos de guardado de métricas mencionadas en la metodología del Capítulo 4 y funciones que almacenen los gráficos de estos datos.

    \end{itemize}

    \item \textbf{WP$_2$: Experimentación, análisis de resultados e interpretación}

    \begin{itemize}

        \item \textbf{A$_{2.1}$} Aplicar metodología de diseño experimental para unificar los datos almacenados post-ejecución, verificando su correctitud.

        \item \textbf{A$_{2.2}$} Ejecutar una larga fase de experimentación con los algoritmos con la propuesta incorporada y la versión original de estos en las instancias seleccionadas de los problemas de optimización combinatorial benchmarks y el Profitable Tour Problem. Conjuntamente ir verificando salidas para evitar errores.

        \item \textbf{A$_{2.3}$} Análisis estadístico comparativo entre el algoritmo con la nueva propuesta y su versión tradicional. Se evalúa la calidad de la solución y su rendimiento a través de la convergencia. Plasmados en tablas y gráficos.

        \item \textbf{A$_{2.4}$} Interpretación de los resultados, conclusiones, posibles consideraciones y trabajos a futuro.

    \end{itemize}

\end{itemize}




%%%%%%%%%% en el ámbito de la hibridación de metaheurísticas y métodos de inteligencia artificial, definición del problema, propuesta de solución, métodos investigativos para ejecutar el desarrollo experimental y la planificación
